\documentclass[11pt]{article}

\usepackage[margin=1in]{geometry}
\usepackage{amsmath}
\usepackage{amssymb}
\usepackage{array}
\usepackage{longtable}

\sloppy

\title{CSE 753 -- Advanced Reinforcement Learning \\ Study Notes: Lecture 09, Reward Learning}
\author{}
\date{}

\begin{document}
\maketitle

\noindent This note covers Lecture 09 (Reward Learning: goal classifiers, learning
reward from human preferences, RLHF, DPO, and RLAIF). It is built from the lecture
slides, the exam-focus guide, and the ``what is expected'' guide. Each concept is
explained twice: once in plain, everyday language, and once with the precise
technical statement a grader expects. Every equation is broken into a symbol table.
Places where the slides contain a typo or an ambiguity are flagged explicitly,
because the exam is open-book and rewards students who notice these and state a
clean assumption.

\bigskip
\noindent\textbf{The one idea that ties this whole lecture together:} whenever we
replace the true objective with a learned stand-in (a classifier, a reward model),
the optimizer will chase the places where the stand-in is \emph{wrong}, not just
the places where it is right. Every failure mode in this lecture -- classifier
exploitation, reward hacking under RLHF -- is this same idea wearing a different
costume. Say this explicitly on the exam whenever it is relevant; it is repeatedly
flagged as a full-marks signal.

\section{Where Does the Reward Come From?}

\subsection{What is the concept?}
In every RL problem we assume a reward function exists and tells the agent how
good each state or action is. This slide asks the question underneath that
assumption: in a real task (a robot, a dialogue agent, a self-driving car), who
actually writes down that reward function?

\subsection{Plain-English explanation}
Think about training a puppy. In a simple game you can say ``+1 point for every
ball fetched.'' But how do you give a robot a ``score'' for making a good cup of
tea, or give a chatbot a ``score'' for giving a helpful answer? Nobody can write a
formula for that. In real tasks we usually only have a hand-made stand-in
(a \emph{proxy}) for what we actually want, because the true goal is too fuzzy to
write as a formula.

\subsection{Why not just imitate an expert instead?}
The slide raises direct imitation learning (copy what an expert did) as one
alternative, and lists why it is not enough by itself:
\begin{itemize}
  \item Imitation only mimics actions; it never reasons about \emph{why} those
  actions were good, so it cannot generalize to a new situation the expert never
  faced.
  \item The expert might move in ways the robot physically cannot (different
  ``degrees of freedom''), so copying literally does not transfer.
  \item Sometimes there is no expert at all to record demonstrations from.
\end{itemize}
This motivates the rest of the lecture: instead of copying an expert, or hand
writing a reward, can we \emph{learn} a reward function from weaker, easier
signals such as example goal states or human comparisons?

\subsection{Problem it solves}
It solves the ``reward specification'' problem: how to get a usable reward signal
for tasks where no formula for the reward exists.

\subsection{Problem it cannot solve (yet)}
Just raising the question does not solve it. This slide is purely diagnostic; the
rest of the lecture is the set of concrete fixes (goal classifiers, preference
learning, RLHF, RLAIF).

\section{Goal Classifiers}

\subsection{What is the concept?}
A goal classifier is a small binary (yes/no) classifier trained to recognize
``did the task succeed in this state?'' Its output is then \emph{used as the
reward signal} for an RL algorithm.

\subsection{Plain-English explanation}
Suppose the task is ``put the pencil case behind the notebook.'' Instead of
writing a mathematical rule for success, you show the computer a pile of pictures
where the task succeeded (label 1) and a pile where it did not (label 0), and you
train an ordinary image classifier on those pictures. Now, whenever the robot ends
up in some new state, you ask the classifier ``does this look like success?'' and
whatever number it spits out (between 0 and 1) becomes the reward for that state.

\subsection{The algorithm, step by step (slide 3)}
\begin{enumerate}
  \item Collect example states inside the goal set $G$ (success) and outside it
  (failure).
  \item Train a binary classifier $C_\phi(s)$ with inputs $s_i$ and labels
  $\mathbf{1}(s_i \in G)$.
  \item Run RL using the classifier's output as the reward.
\end{enumerate}

\subsection*{Equation / notation table}
\begin{center}
\begin{tabular}{|>{\raggedright\arraybackslash}p{3.2cm}|>{\raggedright\arraybackslash}p{10.3cm}|}
\hline
\textbf{Symbol} & \textbf{What it means} \\
\hline
$s_i$ & One example state (a snapshot of the world, e.g. a photo of the table). \\
\hline
$G$ & The set of all states that count as ``the task is done.'' \\
\hline
$\mathbf{1}(s_i \in G)$ & A label that is $1$ if state $s_i$ is a success and $0$
otherwise -- this is exactly what you would type into a spreadsheet next to each
photo. \\
\hline
$C_\phi(s)$ & The trained classifier, with learnable parameters $\phi$; given a
state it outputs a number between 0 and 1. \\
\hline
\end{tabular}
\end{center}
\textbf{What the whole idea is saying:} ``train an ordinary yes/no classifier on
labelled snapshots, and reuse its confidence score as if it were a reward.''

\subsection{What can go wrong (slide 4) -- Attribute the failure}
\textbf{Name of the phenomenon:} the RL agent \emph{exploits the classifier's
weaknesses}.

\textbf{Mechanism, step by step:}
\begin{enumerate}
  \item The RL agent is not told ``reach success states''; it is told ``maximize
  the classifier's output number.''
  \item Those are only the same thing on states the classifier was actually
  trained on.
  \item RL exploration is very good at finding unusual, never-before-seen states.
  \item On those unseen states the classifier's output is meaningless (it is
  guessing), and RL will happily steer toward whichever unseen state happens to
  score high, even though it is not a real success.
\end{enumerate}
\textbf{Canonical example (from the exam-focus guide):} a cup left upside down on
a coaster scores $0.91$ from the classifier even though the task was never
completed -- the classifier had never seen an upside-down cup during training, so
it guessed, and it guessed wrong in a way that happened to reward the agent.

\subsection{The proposed fix (slides 5--6): retrain on visited states}
\textbf{Idea:} every state the RL agent visits during rollouts is added to the
\emph{negative} set $D_-$, the classifier is retrained on a freshly rebalanced
50/50 mix of $D_+$ and $D_-$, and the loop repeats.

\begin{center}
\begin{tabular}{|>{\raggedright\arraybackslash}p{3.2cm}|>{\raggedright\arraybackslash}p{10.3cm}|}
\hline
\textbf{Symbol} & \textbf{What it means} \\
\hline
$D_+$ & The growing set of states known to be successes (goal states). \\
\hline
$D_-$ & The growing set of states known to be failures, continually topped up with
every state the RL agent has ever visited. \\
\hline
$s_t$ & A state visited at time $t$ during a rollout of the current policy $\pi$. \\
\hline
$D_- \leftarrow D_- \cup \{s_t\}$ & ``Add this newly visited state to the failure
pile,'' regardless of whether it happens to actually be a success. \\
\hline
\end{tabular}
\end{center}
\textbf{Plain-English reading of the loop:} every time the agent finds a new
``loophole'' state that fools the classifier, that exact state gets thrown into
the failure pile and the classifier is retrained. The classifier can never be
fooled twice by the same trick, because the trick becomes a labelled negative
example the moment it is used.

\subsection{The paradox and its resolution -- Justify the design (E4)}
\textbf{The worry:} genuine goal states, once visited during rollouts, also get
added to $D_-$. Doesn't that teach the classifier that the goal is a failure?

\textbf{The resolution:} no, because of the 50/50 \emph{rebalancing} step, not
because the goal states ``look different.'' Work through it:
\begin{itemize}
  \item A true goal state can appear on \emph{both} sides: in $D_+$ (from the
  original positive examples, or from being reached during RL and recorded as
  success) and in $D_-$ (from the blanket rule that every visited state is added
  to the negatives).
  \item $D_+$ starts small; $D_-$ grows huge because \emph{every} visited state
  goes there, success or not.
  \item Balancing the training batch to 50/50 means each example in the small
  $D_+$ is effectively \emph{upweighted} relative to each example in the huge
  $D_-$.
  \item For a state that appears on both sides, the upweighted positive evidence
  is at least as strong as the negative evidence, so the best classifier output
  for it is $p \ge 0.5$.
  \item A true failure state, by contrast, appears \emph{only} in $D_-$ and gets
  driven toward $0$.
\end{itemize}
\textbf{Why this matters for exam marks:} an answer that says ``it still works
because the goal state looks different from other states'' earns zero credit.
The question is specifically testing whether you know it is the \emph{balancing
step} doing the work, not anything about the state's appearance.

\subsection{Advantages}
\begin{itemize}
  \item No reward function has to be hand-engineered; only labelled example
  states are needed, which are usually far easier to collect than a formula.
  \item Works for tasks that are easy to \emph{recognize} but hard to
  \emph{describe} mathematically (e.g. ``the room looks tidy'').
\end{itemize}

\subsection{Disadvantages}
\begin{itemize}
  \item Vulnerable to exploitation by the very optimizer it is meant to guide,
  unless the retraining loop (slides 5--6) is used.
  \item Binary only: it can say ``success'' or ``not success'' but cannot rank two
  different successes against each other (a perfectly folded shirt and a barely
  acceptable one score the same).
  \item The retraining loop adds engineering complexity (a classifier retraining
  step nested inside the RL loop) and still needs an initial, reasonably diverse
  $D_+$/$D_-$ to start from.
\end{itemize}

\subsection{What problem it solves}
Turns a task that has no formula for success into a supervised-learning problem:
label some states, train a classifier, use the classifier's confidence as reward.

\subsection{What it cannot solve}
It cannot express \emph{degrees} of success -- there is no way for a goal
classifier to say ``this outcome is better than that outcome,'' only ``this is a
success/failure.'' It also cannot originate the idea of the goal on its own; a
human still has to supply the initial labelled examples.

\subsection{What problem it gave rise to}
Its binary, all-or-nothing nature is exactly the gap that motivates the next
section: learning reward from \emph{human preferences} (which trajectory is
better?), because preferences give a graded signal instead of a yes/no stamp.

\subsection{When to use it}
Use a goal classifier when success is visually or perceptually easy to recognize
but has no clean symbolic description (most manipulation and navigation tasks),
and when you are able to run the retrain-on-visited-states loop to keep it honest.

\subsection{When not to use it}
Avoid it when you need to rank the \emph{quality} of successes, not just detect
success versus failure (e.g. comparing two different helpful chatbot answers) --
use preference-based reward learning (Section~3) instead.

\subsection{Counterfactual -- what if a part were done differently?}
If step 5 (adding visited states to $D_-$) were skipped, the classifier would
never see the agent's newly-discovered exploits, and the ``what can go wrong''
failure of Section 2.4 would recur indefinitely -- the agent would keep finding
new blind spots forever. If the 50/50 rebalancing were dropped and the raw
(highly imbalanced) $D_+ / D_-$ were used instead, the massive $D_-$ set would
swamp the training signal and the classifier could be pushed to output near-zero
for genuine goal states, since a small number of positive labels for that state
would no longer be able to outweigh a large number of negative labels for it --
which reintroduces exactly the same problem the retraining loop was built to
avoid.

\section{Learning Reward from Human Preferences}

\subsection{What is the concept?}
Instead of asking a human to grade a single trajectory with a number, ask a much
easier question: ``Given two trajectories, which one is better?'' Then fit a
reward function so that it assigns a higher score to whichever one the human
preferred.

\subsection{Plain-English explanation}
It is much harder for a person to say ``this trajectory is worth exactly 7.3
points'' than to say ``I like trajectory A more than trajectory B.'' Comparisons
are a natural, easy judgment for humans; absolute scores are not. So the lecture
builds a reward function purely from these easier, relative comparisons.

\subsection{The reward-consistency requirement (slide 8)}
If a human says $\tau_w$ (the winner) is better than $\tau_l$ (the loser), we want
a reward function $r_\theta$ such that the total reward it assigns to $\tau_w$ is
higher than the total reward it assigns to $\tau_l$:
\[
\sum_{(s,a)\in\tau_w} r_\theta(s,a) \;>\; \sum_{(s,a)\in\tau_l} r_\theta(s,a)
\]

\paragraph{Error flag.} The lecture slide (slide 8) actually writes the
right-hand side as $\sum_{(s,a)\in\tau_w} l_\theta(s,a)$ -- both the sum and the
subscript on the summation are copied from the left-hand side by mistake. The
corrected statement, matching the surrounding text and every later slide, is the
one shown above: the second sum ranges over $\tau_l$ (the losing trajectory), not
$\tau_w$. Write the corrected form on the exam and note in one sentence that the
slide has a typo.

\begin{center}
\begin{tabular}{|>{\raggedright\arraybackslash}p{3.2cm}|>{\raggedright\arraybackslash}p{10.3cm}|}
\hline
\textbf{Symbol} & \textbf{What it means} \\
\hline
$\tau_w$ & The trajectory the human judged to be the \emph{winner} (better one). \\
\hline
$\tau_l$ & The trajectory the human judged to be the \emph{loser} (worse one). \\
\hline
$(s,a)$ & One state-action pair visited along a trajectory. \\
\hline
$r_\theta(s,a)$ & The learned reward model's score for taking action $a$ in state
$s$, with learnable parameters $\theta$. \\
\hline
\end{tabular}
\end{center}
\textbf{Plain reading:} ``whatever reward function we learn, it had better give a
higher total score to the trajectory the human liked more.''

\subsection{The Bradley--Terry preference model}
To turn ``$\tau_w$ should score higher'' into something we can train with
gradient descent, the lecture defines a \emph{probability} of preference using
the logistic sigmoid $\sigma$:
\[
P(\tau_a \succ \tau_b) = \sigma\big(r_\theta(\tau_a) - r_\theta(\tau_b)\big)
\]
and then trains $\theta$ by maximizing the log-probability of the human's actual
choices:
\[
\max_\theta \; \mathbb{E}_{\tau_w,\tau_l}\Big[\log \sigma\big(r_\theta(\tau_w) -
r_\theta(\tau_l)\big)\Big]
\]

\begin{center}
\begin{tabular}{|>{\raggedright\arraybackslash}p{3.2cm}|>{\raggedright\arraybackslash}p{10.3cm}|}
\hline
\textbf{Symbol} & \textbf{What it means} \\
\hline
$\sigma(x) = \frac{1}{1+e^{-x}}$ & The sigmoid function. Turns any real number
into a value between 0 and 1, so it can be read as a probability. \\
\hline
$r_\theta(\tau_a) - r_\theta(\tau_b)$ & How much more reward the model currently
thinks trajectory $a$ deserves compared to trajectory $b$. \\
\hline
$P(\tau_a \succ \tau_b)$ & The model's guess at the probability that a human would
say trajectory $a$ is better than trajectory $b$. \\
\hline
$\mathbb{E}_{\tau_w,\tau_l}[\cdot]$ & Average over many human-labelled comparison
pairs in the dataset. \\
\hline
$\log \sigma(\cdot)$ & The log-likelihood of the human's actual answer; maximizing
this pushes the reward gap in the direction the human indicated. \\
\hline
\end{tabular}
\end{center}
\textbf{Plain-English reading of the whole equation:} ``the bigger the reward gap
the model assigns between the winner and the loser, the closer $\sigma$ gets to
1, i.e. the more confidently the model agrees with the human; train the model to
make this agreement as strong as possible, on average, across all the
comparisons we have.''

\subsection{Worked numeric example (Execute -- E1)}
Given $r_\theta(\tau_w) = 4.0$ and $r_\theta(\tau_l) = 1.5$:
\begin{align*}
\Delta &= r_\theta(\tau_w) - r_\theta(\tau_l) = 4.0 - 1.5 = 2.5 \\
\sigma(2.5) &= \frac{1}{1+e^{-2.5}} \approx 0.924 \\
\text{loss} &= -\log(0.924) \approx 0.079
\end{align*}
\textbf{Reading the number (E2):} a small loss like $0.079$ means the reward
model already ranks this pair the way the human did -- there is not much left to
learn from this particular comparison. A large loss would mean the model
currently disagrees with the human and this pair carries a strong training
signal.

\subsection{The reward-learning training loop (slide 9)}
\begin{enumerate}
  \item From a dataset of trajectories, sample a batch of $k$ trajectories and
  ask humans to rank them (for LLMs, all $k$ trajectories answer the same
  prompt).
  \item Compute $r_\theta(\tau_1), \dots, r_\theta(\tau_k)$ under the current
  reward model.
  \item For all $\binom{k}{2}$ pairs in the batch, compute the gradient of
  $\mathbb{E}_{\tau_w,\tau_l}[\log \sigma(r_\theta(\tau_w)-r_\theta(\tau_l))]$.
  \item Update $\theta$ using the computed gradient.
  \item Repeat from step 2.
\end{enumerate}
\textbf{Combinatorics (Execute):} ranking $k$ items produces $\binom{k}{2}$
pairwise comparisons. For $k=6$: $\binom{6}{2} = \frac{6 \times 5}{2} = 15$
comparisons. This is why ranking a whole batch at once is more data-efficient
than asking for one pairwise comparison at a time -- a single ranking of 6 items
is worth 15 individual comparisons for free.

\subsection{Why preferences instead of absolute scores -- Justify the design}
\begin{itemize}
  \item Relative comparisons are a far easier, more natural judgment for a human
  to make than an absolute numeric score.
  \item The reward model only needs to be \emph{discriminative} (correctly order
  pairs), not \emph{calibrated} (produce a score on any particular numeric
  scale) -- a much weaker and easier requirement to satisfy.
\end{itemize}

\subsection{Advantages}
\begin{itemize}
  \item Sidesteps the impossible task of hand-writing a numeric reward.
  \item Cheap, natural signal to collect from ordinary humans (or, later, from
  another AI -- Section 6).
  \item Ranking a batch of $k$ items yields $\binom{k}{2}$ comparisons ``for
  free,'' so it is data-efficient.
\end{itemize}

\subsection{Disadvantages}
\begin{itemize}
  \item Only captures the \emph{difference} between two rewards -- the absolute
  scale of the reward function is never pinned down by this loss (this same fact
  reappears, with larger consequences, in the DPO derivation in Section 5).
  \item Needs many human comparisons, which is still a real cost at scale, and
  motivated Section 6 (RLAIF) as a way to remove that cost.
\end{itemize}

\subsection{What problem it solves}
Produces a trainable, differentiable reward function from human judgments that
never required a single absolute numeric score.

\subsection{What it cannot solve}
By itself, this only trains a reward model; it does not yet say how to use that
reward model to actually improve a policy. That is the job of the RLHF pipeline
in the next sections.

\subsection{What problem it gave rise to}
Once you have a learned reward model instead of a true reward, you have created a
new failure mode: the policy can be pushed to \emph{over-optimize} the learned
model rather than the real objective it stands in for. This is the reward-hacking
problem addressed in Section 5.

\subsection{Counterfactual}
If the loss only used $r_\theta(\tau_w)$ alone (with no comparison to
$r_\theta(\tau_l)$), the model would have no notion of relative quality at all --
it would degenerate into trying to make every trajectory's reward large, which is
not a meaningful training signal, since there would be nothing to push the reward
of bad trajectories down.

\section{RLHF: The Three-Stage Pipeline for LLMs}

\subsection{What is the concept?}
For large language models, reward learning from preferences is embedded in a
larger three-stage recipe: pre-train, then supervised fine-tune, then apply
reinforcement learning against a learned reward model. This is the standard
``RLHF'' pipeline.

\begin{center}
\begin{longtable}{|>{\raggedright\arraybackslash}p{3.2cm}|>{\raggedright\arraybackslash}p{4.8cm}|>{\raggedright\arraybackslash}p{5.3cm}|}
\hline
\textbf{Stage} & \textbf{Objective} & \textbf{Data / signal} \\
\hline
1. Large-scale pre-training & Next-token prediction & Massive, mixed-quality,
unlabelled text corpus \\
\hline
2. Supervised fine-tuning & Maximum likelihood on demonstrations & Curated,
higher-quality (prompt, response) pairs \\
\hline
3. RL from human feedback & Maximize the learned reward model minus a KL penalty
to the pre-trained model & Human \emph{pairwise preference comparisons}
$(x, y_w, y_l)$, used to train the reward model \\
\hline
\end{longtable}
\end{center}

\paragraph{Error flag.} On the lecture slide, stage 3's subtitle literally repeats
stage 2's subtitle (``on higher-quality (prompt, response)''). This is a
copy-paste typo. Stage 3's actual data is \emph{preference comparisons}, not
plain (prompt, response) pairs -- do not copy the slide's subtitle verbatim on
the exam.

\subsection{Plain-English explanation of the pipeline}
\begin{enumerate}
  \item \textbf{Pre-training:} the model first reads huge amounts of text and
  just learns to predict the next word. It becomes fluent, but it has not yet
  learned to be a good assistant.
  \item \textbf{Supervised fine-tuning:} the model is shown a smaller set of
  high-quality example conversations and trained to imitate them directly, the
  same way ordinary supervised learning works.
  \item \textbf{RLHF:} finally, the model generates several candidate answers to
  a prompt, a reward model (trained from human preference comparisons, as in
  Section 3) scores them, and reinforcement learning nudges the model's
  parameters toward answers that score higher.
\end{enumerate}

\subsection{Gathering preference data for LLMs (slides 10--11)}
For a prompt $x$ (e.g. ``How do I cook an omelette?''), the model samples two
candidate replies $y$ and $y'$. A human is shown both and picks the better one,
written $y \succ y'$. A reward model $r(x,y)$ (or $RM_\phi(x,y)$) is then trained
on many such comparisons, exactly following the Bradley--Terry loss from
Section 3.

\subsection{Advantages of the three-stage pipeline}
\begin{itemize}
  \item Separates concerns cleanly: fluency (stage 1), instruction-following
  style (stage 2), and alignment with nuanced human judgment (stage 3).
  \item Reuses cheap, plentiful data (raw text) for the most expensive-to-train
  stage, and saves the costly human-labelled data for the stages that need it
  most.
\end{itemize}

\subsection{Disadvantages}
\begin{itemize}
  \item Three separate training stages: expensive, slow, and each stage can
  introduce its own errors that propagate forward.
  \item Stage 3 needs a working reward model \emph{and} an RL loop, both of which
  add engineering complexity and instability (see Section 5).
\end{itemize}

\subsection{What problem it solves}
Converts human notions of ``a good response'' -- which have no formula -- into a
concrete training signal a language model can be optimized against, without
needing millions of hand-written demonstrations for every possible prompt.

\subsection{What it cannot solve}
It does not, by itself, prevent the RL stage from over-optimizing the learned
reward model (reward hacking) -- see Section 5.

\subsection{Counterfactual}
If stage 3 skipped the reward model and instead reinforced the model directly
against the raw preference labels (i.e. treated preference as the reward with no
learned generalization), the model could only ever be evaluated on the exact
pairs humans labelled -- it would have no way to score a brand-new response the
policy generates during RL, so RL could not proceed at all. The learned reward
model is precisely what lets the signal generalize to unseen responses.

\section{RLHF: Optimizing the Learned Reward Model}

\subsection{What is the concept?}
Given a frozen reward model $RM_\phi$, RLHF fine-tunes a copy of the language
model, $p^{RL}_\theta$, to produce responses that score highly under $RM_\phi$,
while adding a penalty that stops it from drifting too far from the original
pre-trained model $p^{PT}$.

\subsection{The objective (slide 12)}
\[
\mathbb{E}_{\hat y \sim p^{RL}_\theta(\hat y \mid x)}
\left[RM_\phi(x,\hat y) - \beta \log\left(\frac{p^{RL}_\theta(\hat y \mid
x)}{p^{PT}(\hat y \mid x)}\right)\right]
\]

\begin{center}
\begin{longtable}{|>{\raggedright\arraybackslash}p{3.2cm}|>{\raggedright\arraybackslash}p{10.3cm}|}
\hline
\textbf{Symbol} & \textbf{What it means} \\
\hline
$p^{PT}(\hat y \mid x)$ & The frozen, pre-trained (or instruction-tuned) language
model; the fixed ``anchor'' we do not want to drift too far from. \\
\hline
$p^{RL}_\theta(\hat y \mid x)$ & The copy of the model currently being trained by
RL; starts identical to $p^{PT}$, then its parameters $\theta$ are updated. \\
\hline
$\hat y$ & A response generated by the current RL policy for prompt $x$. \\
\hline
$RM_\phi(x,\hat y)$ & The learned reward model's score for how good response
$\hat y$ is for prompt $x$. \\
\hline
$\log\left(\dfrac{p^{RL}_\theta(\hat y|x)}{p^{PT}(\hat y|x)}\right)$ & A per-response
measure of how much more (or less) likely the RL model makes this response
compared to the original pre-trained model -- large when the RL model has drifted
away from the original behaviour. \\
\hline
$\beta$ & A single number (a ``price tag'') controlling how strongly drift away
from $p^{PT}$ is penalized. \\
\hline
\end{longtable}
\end{center}
\textbf{Plain-English reading:} ``try to make the model's answers score highly on
the learned reward model, but charge a penalty, priced at $\beta$, for every bit
the model's behaviour drifts away from how it used to talk before RL started.''

\subsection{Reading the two extremes of $\beta$ -- Bound the claim (E5)}
\begin{itemize}
  \item \textbf{$\beta \to 0$:} the drift penalty disappears; the objective
  becomes pure reward maximization. The policy is then free to exploit any error
  in $RM_\phi$, producing degenerate, high-scoring text that has lost fluency and
  general capability.
  \item \textbf{$\beta \to \infty$:} the penalty completely dominates the reward
  term, forcing $p^{RL}_\theta \approx p^{PT}$. The model stops changing at all --
  it is maximally safe, but it has learned nothing from the human preferences.
\end{itemize}
This is exactly the same shape of trade-off as the KL-constraint limits in
actor-critic methods (Lecture 07): a constraint that is too loose lets the
optimizer run wild against an imperfect target; a constraint that is too tight
prevents any learning.

\subsection{Reward hacking / over-optimization -- Attribute the failure (E3)}
\textbf{Name of the phenomenon:} reward hacking, also called over-optimization of
a learned reward model.

\textbf{Mechanism, step by step:}
\begin{enumerate}
  \item $RM_\phi$ was trained on a finite sample of human comparisons, so it is
  only accurate \emph{on the distribution of responses it was trained on}.
  \item The RL policy is explicitly searching for the responses that score
  highest under $RM_\phi$.
  \item That search process actively drives the policy toward unusual,
  off-distribution responses, precisely where $RM_\phi$'s accuracy has never been
  checked.
  \item The result is a response that $RM_\phi$ rates very highly but that a
  human would actually judge as low quality, confidently wrong, or strange.
\end{enumerate}
\textbf{Diagnostic signs to look for (two together, not one alone):}
\begin{enumerate}
  \item The reward-model score keeps climbing while human-judged quality
  actually falls (confident but wrong outputs).
  \item Outputs start looking qualitatively different from the pre-trained
  model's style, i.e. the KL divergence from $p^{PT}$ has grown large, often with
  repetitive or degenerate stylistic tics.
\end{enumerate}
\textbf{Fix:} increase $\beta$, and/or stop training earlier, and/or retrain the
reward model on fresh, on-policy comparisons so it stays accurate on the region
the policy currently occupies.

\subsection{The shared-structure connection -- worth stating explicitly}
This is the same failure, in a new outfit, as classifier exploitation in
Section 2: a learned stand-in for the true objective is only reliable on the data
it was trained on, and an optimizer that searches hard enough will always find
the stand-in's blind spots. On the exam, naming this connection explicitly (in
one sentence) is called out as a signal of real understanding rather than rote
recall.

\subsection{Advantages}
\begin{itemize}
  \item Gives a tunable dial ($\beta$) that trades off ``learn from preferences''
  against ``stay safe and close to the original model.''
  \item Works even though the true, ideal reward is unknown -- it only needs a
  learned proxy.
\end{itemize}

\subsection{Disadvantages}
\begin{itemize}
  \item Requires careful tuning of $\beta$; get it wrong in either direction and
  you get degenerate outputs or no learning at all.
  \item Needs a full RL training loop plus sampling from the policy at every
  step -- computationally expensive and can be unstable.
  \item Vulnerable to reward hacking whenever $\beta$ is too small or training
  runs too long.
\end{itemize}

\subsection{What problem it solves}
Gives a concrete, tunable way to fine-tune a language model against a learned
notion of ``good response'' without completely destroying its original fluency
and capabilities.

\subsection{What it cannot solve}
It cannot guarantee the learned reward model stays accurate forever as the policy
moves; nothing in the objective itself detects when $RM_\phi$ has become
unreliable -- that has to be watched for externally (the two diagnostic signs
above).

\subsection{What problem it gave rise to}
The expense and instability of running a full RL loop against a separately
trained, frozen reward model is exactly what motivates Direct Preference
Optimization (Section 5): can we skip the reward model and the RL loop entirely?

\subsection{When to use it}
Use full RLHF with an explicit reward model and RL loop when you want the
flexibility to reuse the reward model for multiple downstream policies, or when
online sampling during training is feasible and affordable.

\subsection{When not to use it}
Avoid it when training compute or stability is limited, or when a simpler,
directly-supervised alternative (DPO) would reach a similar result more cheaply.

\subsection{Counterfactual}
If the KL penalty term were removed entirely and only $RM_\phi(x,\hat y)$ were
maximized, that is exactly the $\beta \to 0$ case above: reward hacking would set
in quickly since nothing anchors the model to sensible, fluent language.

\section{Direct Preference Optimization (DPO)}

\subsection{What is the concept?}
DPO is a way to get the same effect as the full RLHF pipeline (fine-tune a
language model to match human preferences) \emph{without} training a separate
reward model and \emph{without} running an RL loop at all. It reduces the whole
problem to a single supervised classification-style loss computed directly on
preference pairs.

\subsection{Plain-English explanation}
The big insight is: the RLHF objective from Section 5 has an exact mathematical
solution on paper -- if you already knew the ideal reward model, you could write
down the exact best-fine-tuned model in one formula, with no trial-and-error RL
needed. DPO flips that formula around: instead of using it to compute the best
model from a known reward, it uses it to say ``whatever the reward model would
have been, it can be written directly in terms of the model itself.'' That
means the model becomes its own reward model, and the preference loss from
Section 3 can be trained directly on the language model's parameters -- no
separate reward network, no RL loop, no sampling during training.

\subsection{The derivation, one link at a time (examinable in full)}

\paragraph{Step 1 -- the closed-form optimum (slide 13).}
The RLHF objective from Section 5 has an exact best solution:
\[
p^*(\hat y \mid x) = \frac{1}{Z(x)}\, p^{PT}(\hat y \mid x)\, \exp\!\left(\frac{1}{\beta}
RM(x,\hat y)\right)
\]

\begin{center}
\begin{longtable}{|>{\raggedright\arraybackslash}p{3.2cm}|>{\raggedright\arraybackslash}p{10.3cm}|}
\hline
\textbf{Symbol} & \textbf{What it means} \\
\hline
$p^*(\hat y \mid x)$ & The exact, ideal fine-tuned model that would come out of
solving the RLHF objective perfectly. \\
\hline
$Z(x)$ & The normalizing constant (partition function) that makes $p^*(\cdot|x)$
sum to 1 over all possible responses $y$; formally
$Z(x) = \sum_y p^{PT}(y|x)\exp(RM(x,y)/\beta)$. It depends only on the prompt
$x$, not on any particular response. \\
\hline
$\exp\!\left(\frac{1}{\beta}RM(x,\hat y)\right)$ & Reshapes the pre-trained
distribution to favour high-reward responses more strongly as $\beta$ shrinks. \\
\hline
\end{longtable}
\end{center}
\textbf{Plain reading:} ``the perfect fine-tuned model just takes the original
model's probabilities and re-weights them to favour high-reward responses,
renormalized so everything still sums to 1.''

\paragraph{Step 2 -- rearrange to express the reward as a log-ratio.}
\[
RM(x,\hat y) = \beta \log\left(\frac{p^*(\hat y \mid x)}{p^{PT}(\hat y \mid
x)}\right) + \beta \log Z(x)
\]
This is just algebra on Step 1: solve for $RM(x,\hat y)$ instead of for
$p^*(\hat y|x)$.

\paragraph{Step 3 -- this holds for an arbitrary language model.}
Since the relationship in Step 2 is just algebra, it holds no matter which model
plays the role of $p^*$. So we can plug in \emph{any} candidate model
$p^{RL}_\theta$ and \emph{define} an implied reward model straight from it:
\[
RM_\theta(x,\hat y) = \beta \log\left(\frac{p^{RL}_\theta(\hat y \mid
x)}{p^{PT}(\hat y \mid x)}\right) + \beta \log Z(x)
\]
\textbf{Plain reading:} ``the policy is implicitly its own reward model'' -- once
we know $\beta$, the pre-trained model, and the current policy, we already know
what reward model would have produced this policy as the RLHF optimum. No
separate reward network is needed.

\paragraph{Step 4 -- substitute into the Bradley--Terry loss; $Z(x)$ cancels.}
\[
RM_\theta(x,y^w) - RM_\theta(x,y^l) = \beta \log \frac{p^{RL}_\theta(y^w \mid
x)}{p^{PT}(y^w \mid x)} - \beta \log \frac{p^{RL}_\theta(y^l \mid x)}{p^{PT}(y^l
\mid x)}
\]
\[
J_{\text{DPO}}(\theta) = -\mathbb{E}_{(x,y^w,y^l)\sim\mathcal{D}}\left[\log
\sigma\Big(RM_\theta(x,y^w) - RM_\theta(x,y^l)\Big)\right]
\]

\begin{center}
\begin{longtable}{|>{\raggedright\arraybackslash}p{3.2cm}|>{\raggedright\arraybackslash}p{10.3cm}|}
\hline
\textbf{Symbol} & \textbf{What it means} \\
\hline
$y^w$, $y^l$ & The human-preferred (winner) and non-preferred (loser) response to
the same prompt $x$. \\
\hline
$\mathcal{D}$ & The dataset of human preference triples $(x, y^w, y^l)$. \\
\hline
$J_{\text{DPO}}(\theta)$ & The final DPO training loss -- an ordinary supervised
loss, computed directly on the language model's own output probabilities, with no
RL loop and no separate reward network. \\
\hline
\end{longtable}
\end{center}

\subsection{Why $Z(x)$ cancels -- Justify the design (E4)}
\textbf{What $Z(x)$ is:} the partition function for prompt $x$: a sum over the
\emph{entire space} of possible responses $y$. That space is exponentially large
(every possible sequence of tokens), so $Z(x)$ is intractable to compute exactly.

\textbf{Why it cancels:} $Z(x)$ depends only on the prompt $x$, not on the
particular response $y$. The DPO loss only ever needs the \emph{difference}
between the implied reward of two responses to the \emph{same} prompt $x$
($y^w$ and $y^l$). Since $\beta\log Z(x)$ is added identically to both
$RM_\theta(x,y^w)$ and $RM_\theta(x,y^l)$, it is subtracted from itself in the
difference and disappears -- without ever having to be computed.

\textbf{The practical consequence, the part most often missed:} because only the
difference survives, DPO training can compute \emph{relative} rewards between two
responses to the same prompt exactly, but the \emph{absolute} reward of a single
response is only known up to the unknown constant $\beta \log Z(x)$. That is
completely sufficient for preference optimization (which only ever compares
pairs), but it is \emph{not} sufficient for any downstream use that needs a
calibrated, absolute reward score for a single response on its own.

\subsection{Worked numeric example (Execute -- E1)}
Given $\beta = 0.5$ and log-ratio values $\beta$-scaled to $1.2$ (for $y^w$) and
$0.2$ (for $y^l$):
\begin{align*}
\text{argument} &= 0.5 \times (1.2 - 0.2) = 0.5 \\
\sigma(0.5) &= \frac{1}{1+e^{-0.5}} \approx 0.622 \\
J_{\text{DPO}} &= -\log(0.622) \approx 0.475
\end{align*}
\textbf{Reading the number:} the loss is moderate-sized, meaning the model
currently favours $y^w$ over $y^l$ by only a small margin -- there is meaningful
training signal left in this pair, unlike the earlier Bradley--Terry example
where the reward gap was already large and confident.

\subsection{DPO's advantage over full RLHF, in one sentence}
DPO removes the separate reward-model training stage and the entire RL sampling
loop, turning preference alignment into a single supervised, classification-style
loss computed directly on log-probabilities -- with no reward hacking of a
separate frozen reward model, and generally more stable and cheaper to train.

\subsection{Advantages}
\begin{itemize}
  \item No separate reward network to train, store, or maintain.
  \item No RL loop and no on-policy sampling during training -- ordinary
  supervised-style optimization.
  \item Cannot hack a reward model that does not exist as a separate object.
\end{itemize}

\subsection{Disadvantages}
\begin{itemize}
  \item Only ever learns \emph{relative} preferences between pairs; it never
  produces a calibrated absolute reward for a single response.
  \item Still fundamentally needs human (or AI) preference-pair data -- it does
  not remove the labelling requirement, only the reward-model and RL-loop
  machinery built on top of it.
  \item Because there is no explicit reward model, tools that rely on an
  explicit, standalone reward signal (e.g. for reward shaping or auditing) are
  not directly available.
\end{itemize}

\subsection{What problem it solves}
Removes the cost, complexity, and reward-hacking risk of the separate reward
model plus RL loop in standard RLHF, while producing (in principle) the same
optimal fine-tuned model.

\subsection{What it cannot solve}
Because absolute reward is unrecoverable (only known up to $\beta\log Z(x)$), DPO
cannot supply a calibrated reward score for any single response in isolation --
if a downstream system genuinely needs an absolute quality score (not just
better-or-worse comparisons), DPO is the wrong tool and an explicit reward model
is still required.

\subsection{What problem it gave rise to}
DPO still depends on a fixed, pre-collected preference dataset; it has no
built-in mechanism to acquire fresh preference labels for new, off-distribution
responses the way an online RLHF loop with a live reward model can. If the
policy drifts far from the data the preference pairs were collected on, DPO has
no online correction (unlike RLHF, which the note on slide 9 mentions can be run
inside the loop of online RL).

\subsection{When to use it}
Use DPO when a fixed preference dataset is already available and the goal is a
stable, cheap, single-stage fine-tune, with no need for an explicit, standalone
reward model afterward.

\subsection{When not to use it}
Avoid DPO when the application specifically needs an absolute, calibrated reward
score (e.g. to combine with other scoring systems, or for reward shaping in a
downstream RL problem), or when preference data needs to be collected
continually, on-policy, as training proceeds.

\subsection{Counterfactual -- what if a part were done differently?}
If Step 3 had substituted $p^{RL}_\theta$ for a response to a \emph{different}
prompt $x'$ instead of the same prompt $x$, the $Z(x)$ terms would not match
($\beta\log Z(x) \ne \beta\log Z(x')$) and would not cancel -- this is precisely
why DPO's preference pairs must always compare two responses to the \emph{same}
prompt.

\section{Learning Reward from AI Feedback (RLAIF / Constitutional RL)}

\subsection{What is the concept?}
Instead of asking a human ``which response is less harmful?'', ask another
language model the same question, and use its answer in place of a human
preference label everywhere in the pipeline above (Sections 3--5).

\subsection{The key insight (slide 15), quoted as the exam wants it}
\emph{``Critique is easier than generation.''} A model that cannot reliably
\emph{produce} the best possible response can still reliably \emph{judge} which
of two given responses is less harmful. That asymmetry between generating and
judging is the entire justification for letting an AI stand in for a human
annotator.

\subsection{Plain-English explanation}
Imagine a student who cannot write a great essay from scratch, but who can still
tell you, given two draft essays, which one reads better. Judging is often a much
easier skill than creating. RLAIF leans on exactly this: even if a language model
is not good enough to write the best possible answer, it can often be trusted to
say ``answer A is safer than answer B.''

\subsection{Advantages at scale}
\begin{itemize}
  \item Essentially unlimited throughput, at a tiny marginal cost compared to
  paying and scheduling human annotators.
  \item Consistent application of a single, written policy or ``constitution,''
  with no annotator-to-annotator variance or annotator fatigue over long
  sessions.
  \item Scales instantly to new content volumes.
  \item Avoids exposing human moderators to large volumes of harmful or
  disturbing content.
\end{itemize}

\subsection{The unique risk -- Bound the claim (E5)}
\textbf{The one risk that does not apply to human annotators:} the AI
annotator's own biases and blind spots are \emph{systematic and correlated
across every single label it produces}. Independent human annotators make noisy,
largely independent mistakes that tend to average out across a large dataset. A
single AI annotator model does not have that property -- if it has a blind spot,
that exact blind spot shows up identically in every label it touches. The reward
model then inherits that bias wholesale, and the RL policy optimizes directly
against it, with nothing independent left anywhere in the loop to catch the
drift.

\subsection{Result from the slides (slide 15, Pareto plot)}
Constitutional RL (AI feedback with a written constitution) is shown tracing out
a curve strictly above and to the right of standard RLHF on a plot of
helpfulness versus harmlessness -- described on the slide as a \emph{Pareto
improvement}: better on both axes at once, not merely a trade of one for the
other.

\subsection{Comparison table: goal classifiers vs. RLHF}
\begin{center}
\begin{longtable}{|>{\raggedright\arraybackslash}p{3.4cm}|>{\raggedright\arraybackslash}p{5.0cm}|>{\raggedright\arraybackslash}p{5.0cm}|}
\hline
\textbf{Aspect} & \textbf{Goal classifier} & \textbf{RLHF (human preferences)} \\
\hline
Type of human input & A set of labelled example states (success/failure), given
once, up front & Ongoing pairwise comparisons between rollouts or responses \\
\hline
Handles nuanced quality? & No -- binary success/failure only; cannot rank two
successes against each other & Yes -- relative preferences give a graded,
continuous reward signal \\
\hline
Main failure mode & Agent exploits states the classifier was never trained on &
Reward hacking / over-optimization of the learned reward model \\
\hline
\end{longtable}
\end{center}
\textbf{The connection to notice:} these are the \emph{same underlying failure}
(a learned proxy gets exploited by the optimizer built on top of it) showing up
at two different scales of the same idea.

\subsection{Advantages}
\begin{itemize}
  \item Removes the human-labelling bottleneck almost entirely.
  \item More consistent than a pool of human annotators.
\end{itemize}

\subsection{Disadvantages}
\begin{itemize}
  \item Inherits and amplifies the annotator model's own systematic biases,
  correlated across the entire dataset.
  \item No independent ground truth left in the loop to detect when the
  annotator model itself has drifted or is wrong.
\end{itemize}

\subsection{What problem it solves}
Removes the cost and throughput bottleneck of collecting human preference labels
at the scale needed for large models.

\subsection{What it cannot solve}
It cannot supply an independent check on the annotator's own judgment -- if the
annotator model is systematically biased or has a blind spot, RLAIF has no
built-in way to detect or correct that, unlike a diverse pool of independent
human annotators whose individual errors tend to cancel out.

\subsection{What problem it gave rise to}
The systematic-bias risk described above is a new, distinct failure mode that a
purely human-labelled pipeline did not have to the same degree -- it motivates
ongoing research into diversifying or auditing AI annotators, and combining AI
and human feedback rather than relying on either alone.

\subsection{When to use it}
Use RLAIF when human-labelling throughput is the bottleneck, when a clear written
policy (a ``constitution'') can specify what ``better'' means, and when some risk
of correlated annotator bias is acceptable or can be independently audited.

\subsection{When not to use it}
Avoid relying on RLAIF alone when independent ground truth or auditability is
critical, or when the judging model itself is known to share the same blind
spots you are trying to correct for in the policy being trained.

\subsection{Counterfactual}
If the AI annotator were replaced by a diverse panel of several different AI
models (rather than one), the correlated-bias risk described above would be
reduced, because different models are less likely to share the exact same blind
spots -- much closer to how independent human annotators' errors average out.

\section{Summary: The Thread Running Through the Whole Lecture}

\begin{center}
\begin{longtable}{|>{\raggedright\arraybackslash}p{4.2cm}|>{\raggedright\arraybackslash}p{9.3cm}|}
\hline
\textbf{Stage} & \textbf{The recurring pattern} \\
\hline
Hand-designed reward & Impossible to write for most real tasks. \\
\hline
Goal classifier & A learnable stand-in, but the RL agent exploits states the
classifier was never trained on. \\
\hline
Retrain on visited states & Removes the exploit, but only gives a binary
success/failure signal. \\
\hline
Human preferences (Bradley--Terry) & Gives a graded signal instead of binary, but
needs a full RLHF pipeline (reward model + RL loop) to use it. \\
\hline
RLHF with KL penalty & Works, but the RL policy can over-optimize
(``reward-hack'') the learned reward model, off its valid distribution. \\
\hline
DPO & Removes the separate reward model and RL loop entirely, but only ever
recovers \emph{relative}, not absolute, reward. \\
\hline
RLAIF & Removes the human-labelling bottleneck, but introduces a new, systematic
(not averaging-out) annotator-bias risk. \\
\hline
\end{longtable}
\end{center}

\noindent\textbf{The single sentence to remember:} optimizing a learned proxy for
an objective makes the optimizer seek out exactly the places where the proxy is
wrong -- this is the same phenomenon behind classifier exploitation, RLHF reward
hacking, and (per the exam-focus guide) OOD value overestimation in offline RL
from Lecture 08. Naming this connection, in one sentence, whenever a question
touches any one of these three ideas is repeatedly flagged as the mark of real
understanding rather than memorized recall.

\end{document}